\documentclass[final,5p,times,twocolumn]{elsarticle}

\usepackage{amsmath,amssymb,amsthm,mathtools}
\usepackage{enumitem}
\usepackage{microtype}

\journal{Applied Mathematics Letters}
\newtheorem{theorem}{Theorem}
\newtheorem{lemma}[theorem]{Lemma}
\newtheorem{corollary}[theorem]{Corollary}
\newtheorem{proposition}[theorem]{Proposition}
\newcommand{\Om}{\Omega}
\newcommand{\R}{\mathbb R}
\newcommand{\norm}[1]{\lVert#1\rVert}
\newcommand{\ubar}{\overline u_0}

\begin{document}
\begin{frontmatter}

\title{Mass-weighted damping and a rate dichotomy for chemotaxis with signal-dependent motility}
\author{}

\begin{abstract}
We consider
\[
u_t=\Delta(\varphi(v)u),\qquad v_t=\Delta v+uF(v)
\]
on a smooth bounded domain, with homogeneous Neumann boundary conditions. Positive conserved cell mass yields a coercive mechanism even without a pointwise lower bound for $u$. For every $q\ge2$ we prove a nonlinear mass-weighted Poincar\'e estimate converting $\int u|v-v_*|^q$ into control of the signal energy. Hence quadratic dissipativity gives exponential signal decay, whereas $(s-v_*)F(s)\le-\beta|s-v_*|^{2+\theta}$ yields the algebraic rate $t^{-1/\theta}$. If $u$ is eventually bounded in one finite $L^p$ space with $p>\max\{n,2\}$, the decay propagates to $L^\infty\times W^{1,\infty}$. In the degenerate case every
\[
0<\mu<\frac1\theta\min\left\{1,\frac{2(p-n)}{pn}\right\}
\]
is admissible for the full system. The framework applies to the production--consumption model of Qin and Zheng and to superlinear signal consumption; spatially homogeneous solutions show that the signal exponent $1/\theta$ is optimal.
\end{abstract}

\begin{keyword}
chemotaxis \sep signal-dependent motility \sep weighted damping \sep stabilization \sep algebraic decay
\end{keyword}
\end{frontmatter}

\section{Introduction and main results}
Let $\Om\subset\R^n$, $n\ge1$, be smooth, bounded and connected. We study nonnegative global classical solutions of
\begin{equation}\label{sys}
 u_t=\Delta(\varphi(v)u),\qquad v_t=\Delta v+uF(v),
\end{equation}
with $\partial_\nu u=\partial_\nu v=0$. Set
\[
m:=\int_\Om u_0>0,\qquad \ubar:=m/|\Om|.
\]
The first equation conserves $m$. A recent motivation is the production--consumption law $F(s)=1-\alpha s$ studied by Qin and Zheng \cite{QinZheng2026}, who prove global boundedness under a structural condition on the motility. Stabilization for the same reaction with a classical chemotactic flux was obtained by Tao and Winkler \cite{TaoWinkler2025}; for signal-dependent motility, exponential relaxation is already known for direct consumption $F(s)=-s$ \cite{LiZhao2021}. Our point is therefore not stabilization per se, but a common rate mechanism covering nonzero equilibria, degenerate damping and only eventual finite-$L^p$ control of $u$.

Let $C_P$ satisfy $\norm{f-f_\Om}_2\le C_P\norm{\nabla f}_2$.

\begin{lemma}[Mass-weighted coercivity]\label{lem:coerc}
Let $q\ge2$, $\rho\ge0$, $\int_\Om\rho=m>0$ and $\norm{\rho}_2\le K$. Then
\begin{equation}\label{coerc}
\norm f_2^2\le A\norm{\nabla f}_2^2+B_q\left(\int_\Om\rho|f|^q\right)^{2/q}
\end{equation}
for every $f\in H^1(\Om)$ for which the right-hand side is finite, where
\[
A=C_P^2\left(1+\frac{2|\Om|K^2}{m^2}\right),\qquad B_q=2|\Om|m^{-2/q}.
\]
\end{lemma}

The estimate is elementary; its role is to recover unweighted coercivity from positive mass without requiring $\rho$ to be pointwise positive.

\begin{theorem}[Weighted-damping dichotomy]\label{thm:weighted}
Let $z$ solve
\[
z_t=\Delta z+\rho(x,t)F(z),\qquad \partial_\nu z=0,
\]
on $[T,\infty)$ and remain in a compact interval $I$. Assume $\rho\ge0$, $\int_\Om\rho(\cdot,t)=m>0$, $\sup_{t\ge T}\norm{\rho(t)}_2\le K$, and for some $z_*\in I$, $\beta>0$ and $q\ge2$,
\begin{equation}\label{dissq}
F(z_*)=0,
\qquad (s-z_*)F(s)\le-\beta|s-z_*|^q \quad (s\in I).
\end{equation}
If $q=2$, then $\norm{z(t)-z_*}_2\le Ce^{-\lambda(t-T)}$. If $q>2$, then
\begin{equation}\label{rateweighted}
\norm{z(t)-z_*}_2\le C(1+t-T)^{-1/(q-2)}.
\end{equation}
\end{theorem}

\begin{theorem}[Full stabilization from finite $L^p$ control]\label{thm:full}
Assume that $v(x,t)\in I$ for a compact interval $I$, $\varphi\in C^1(I)$, $\min_I\varphi>0$, $F\in C^1(I)$, and for some finite
\[
p>\max\{n,2\},\qquad \sup_{t\ge T}\norm{u(t)}_p\le U_p.
\]
Let $F(v_*)=0$. If
\[
(s-v_*)F(s)\le-\beta|s-v_*|^2 \quad(s\in I),
\]
then, for $t\ge T+2$,
\begin{equation}\label{full-exp}
\norm{u(t)-\ubar}_\infty+\norm{v(t)-v_*}_{W^{1,\infty}}
\le Ce^{-\lambda(t-T)}.
\end{equation}
If instead, for some $\theta>0$,
\begin{equation}\label{deg}
(s-v_*)F(s)\le-\beta|s-v_*|^{2+\theta},
\end{equation}
then $\norm{v(t)-v_*}_2=O((1+t-T)^{-1/\theta})$, and for every
\begin{equation}\label{murange}
0<\mu<\frac1\theta\min\left\{1,\frac{2(p-n)}{pn}\right\},
\end{equation}
\begin{equation}\label{full-alg}
\norm{u(t)-\ubar}_\infty+\norm{v(t)-v_*}_{W^{1,\infty}}
=O((1+t-T)^{-\mu}).
\end{equation}
No sign condition on $\varphi'$ is used in this stabilization argument.
\end{theorem}

\section{Proofs}
\begin{proof}[Proof of Lemma~\ref{lem:coerc}]
With $f_\Om=|\Om|^{-1}\int_\Om f$, H\"older, Cauchy--Schwarz and Poincar\'e give
\[
m|f_\Om|\le m^{1-1/q}\left(\int_\Om\rho|f|^q\right)^{1/q}+KC_P\norm{\nabla f}_2.
\]
Insert this into the exact identity
$\norm f_2^2=\norm{f-f_\Om}_2^2+|\Om||f_\Om|^2$ and use $(a+b)^2\le2a^2+2b^2$.
\end{proof}

\begin{proof}[Proof of Theorem~\ref{thm:weighted}]
Set $w=z-z_*$, $E=\norm w_2^2$, $G=\norm{\nabla w}_2^2$, and $D=\int_\Om\rho|w|^q$. Testing by $w$ gives
\begin{equation}\label{energy}
\frac12E'+G+\beta D\le0,
\end{equation}
whereas Lemma~\ref{lem:coerc} yields $E\le AG+B_qD^{2/q}$. For $q=2$ this gives a spectral gap and exponential decay. Let $q>2$ and put $X=G+\beta D$. If $X\le1$, then $G\le X$, $D\le X/\beta$ and $X\le X^{2/q}$, hence $E\le CX^{2/q}$. If $X\ge1$, compactness of $I$ gives $E\le C\le CX^{2/q}$. Thus $X\ge cE^{q/2}$ and \eqref{energy} implies $E'\le-2cE^{q/2}$. Integrating $E^{-(q-2)/2}$ proves \eqref{rateweighted}; if $E$ reaches zero, \eqref{energy} keeps it zero.
\end{proof}

\begin{proof}[Proof of Theorem~\ref{thm:full}]
Since $p>2$, the assumed $L^p$ bound gives a uniform $L^2$ bound for $u$, so Theorem~\ref{thm:weighted} yields the stated $L^2$ signal rates. Put $w=v-v_*$. Since $F(v_*)=0$ and $F\in C^1(I)$, $|F(v)|\le L_F|w|$. Choose $n<r<p$ and $s$ by $1/r=1/p+1/s$. Interpolation between the $L^2$ signal rate and the bounded range gives, in the degenerate case,
\begin{equation}\label{wrate}
\norm{w(t)}_s=O((1+t-T)^{-\sigma(r)/\theta}),\quad
\sigma(r):=\min\left\{1,\frac{2(p-r)}{pr}\right\}.
\end{equation}
Hence $\norm{uF(v)}_r$ has the same rate. For $t\ge T+1$,
\[
w(t)=e^{\Delta}w(t-1)+\int_{t-1}^t e^{(t-\tau)\Delta}u(\tau)F(v(\tau))\,d\tau.
\]
The standard Neumann estimate \cite{KeLiWang2022}
\[
\norm{\nabla e^{\tau\Delta}f}_\infty
\le C(1+\tau^{-1/2-n/(2r)})\norm f_r,
\quad0<\tau\le1,
\]
has an integrable singularity because $r>n$. Fixed-time smoothing of the first term and \eqref{wrate} therefore give
\[
\norm{\nabla v(t)}_\infty=O((1+t-T)^{-\sigma(r)/\theta}).
\]
Together with $|w_\Om|\le|\Om|^{-1/2}\norm w_2$ and $\norm{w-w_\Om}_\infty\le C\norm{\nabla w}_\infty$, this gives the same $W^{1,\infty}$ rate. Letting $r\downarrow n$ yields every $\mu$ in \eqref{murange}; the quadratic case is identical with exponential interpolation.

Set $q_u=u-\ubar$. Integration by parts gives
\[
\frac12\frac d{dt}\norm{q_u}_2^2
=-\int_\Om\varphi(v)|\nabla q_u|^2
-\int_\Om u\varphi'(v)\nabla q_u\cdot\nabla v.
\]
Since $1/2=1/p+1/s_0$ with $s_0=2p/(p-2)$, H\"older, Young and Poincar\'e imply
\begin{equation}\label{cellL2}
\frac d{dt}\norm{q_u}_2^2+c\norm{q_u}_2^2
\le C\norm{\nabla v}_\infty^2.
\end{equation}
Thus $q_u$ inherits the same exponential, respectively admissible algebraic, decay in $L^2$.

For the uniform upgrade write
\[
(q_u)_t-\nabla\cdot(a\nabla q_u)=\nabla\cdot B,
\quad a=\varphi(v),\quad B=u\varphi'(v)\nabla v.
\]
The coefficient $a$ is uniformly elliptic, $\norm{B(t)}_p\le C\norm{\nabla v(t)}_\infty$, and the boundary condition is
\[
(a\nabla q_u+B)\cdot\nu
=\varphi(v)\partial_\nu u+u\varphi'(v)\partial_\nu v=0.
\]
Choose $Q_*=4p/(p-n)$; then $Q_*>2$ and $n/p+2/Q_*<1$. If $n\ge2$, Choi's boundary local-boundedness theorem \cite[Theorem~1.1]{Choi2016}, on fixed backward cylinders, gives
\begin{equation}\label{choi}
\norm{q_u}_{L^\infty(Q_{r_0/2})}
\le C\norm{q_u}_{L^2(Q_{r_0})}+C\norm B_{L^{p,Q_*}(Q_{r_0})}.
\end{equation}
Here $r_0>0$ is independent of the terminal time because all lower-order coefficients in Choi's operator vanish. For $n=1$, the same De Giorgi proof applies after replacing its dimension-dependent mixed-norm Sobolev step by the one-dimensional Gagliardo--Nirenberg estimate
\[
\norm h_{L_t^Q L_x^P}\le C\left(\sup_\tau\norm{h(\tau)}_2+\norm{\partial_xh}_{L^2_{x,t}}\right),
\quad \frac1P+\frac2Q=\frac12,
\]
The estimate follows by applying the one-dimensional Gagliardo--Nirenberg inequality at each time and integrating; the interval measure-density condition is automatic. Thus \eqref{choi} also holds in one dimension. Combining \eqref{cellL2}, the decay of $B$, and a finite covering of $\overline\Om$ proves \eqref{full-exp}--\eqref{full-alg}.
\end{proof}

\section{Applications and sharpness}
\begin{corollary}[Production--consumption]\label{cor:qz}
Let $F(s)=1-\alpha s$, $\alpha>0$. Then $v_*=1/\alpha$ and $(s-v_*)F(s)=-\alpha(s-v_*)^2$. Hence every solution satisfying Theorem~\ref{thm:full} converges exponentially to $(\ubar,1/\alpha)$. In particular, Qin and Zheng's uniformly bounded solutions \cite{QinZheng2026} satisfy \eqref{full-exp}: the maximum principle gives
\[
0\le v(x,t)\le\max\{\norm{v_0}_\infty,1/\alpha\}.
\]
Their structural motility assumptions remain relevant for boundedness, but are not used in the subsequent stabilization stage.
\end{corollary}

\begin{corollary}[Superlinear consumption]\label{cor:super}
Under the remaining hypotheses of Theorem~\ref{thm:full}, suppose $I\subset[0,\infty)$ and $F(s)=-s^m$, $m>1$. Then $v_*=0$ and \eqref{deg} holds with $\theta=m-1$. Therefore
\[
\norm{v(t)}_2=O((1+t-T)^{-1/(m-1)}),
\]
and \eqref{full-alg} holds for every
\[
0<\mu<\frac1{m-1}\min\left\{1,\frac{2(p-n)}{pn}\right\}.
\]
\end{corollary}

\begin{proposition}[Sharp signal exponent]\label{prop:sharp}
Let $F(s)=-\kappa(s-v_*)|s-v_*|^\theta$, $\kappa,\theta>0$. For homogeneous data $u_0\equiv\ubar>0$ and $v_0-v_*\equiv w_0\ne0$ whose signal trajectory stays in the domain of $F$,
\[
|v(t)-v_*|=\left(|w_0|^{-\theta}+\theta\kappa\ubar t\right)^{-1/\theta}.
\]
Hence the signal exponent $1/\theta$ cannot in general be improved.
\end{proposition}

\section{Conclusion}
Positive conserved mass converts weighted signal dissipation into a closed energy inequality, independently of pointwise positivity of the cell density. This separates boundedness from asymptotic relaxation and makes the transition from exponential to algebraic rates transparent. The production--consumption model and superlinear consumption arise as the nondegenerate and degenerate branches of the same mechanism.

\bibliographystyle{elsarticle-num}
\bibliography{references}
\end{document}
